\documentclass[a4paper, 12pt]{article}
\pagestyle{empty}
\usepackage[margin=0.5in]{geometry}
\usepackage{graphicx}
\usepackage{amsfonts, amssymb, amsmath, amsthm}
\usepackage{tikz}
\usepackage{minted}

\newcommand{\ceil}[1]{\lceil #1 \rceil}
\newcommand{\floor}[1]{\lfloor #1 \rfloor}
\newcommand{\lt}{\left}
\newcommand{\rt}{\right}
\newcommand{\C}{\mathbb{C}}
\newcommand{\R}{\mathbb{R}}
\newcommand{\Q}{\mathbb{Q}}
\newcommand{\Z}{\mathbb{Z}}
\newcommand{\N}{\mathbb{N}}
\newcommand{\epf}{\hfill \qed}

\begin{document}

\section*{Chapter 1 Classical Algebra}

\subsection*{Exercise 5}

Let $S = \{p+q\alpha+r\alpha^2 \mid p,q,r\in\Q\}$. We first prove that $S$ is a subring of $\C$. 

Note that $1=1+0\alpha+0\alpha^2\in S$. Let $p+q\alpha+r\alpha^2,p'+q'\alpha+r'\alpha^2 \in S$, where $p,q,r,p',q',r' \in \Q$. Then we have $(p+q\alpha+r\alpha^2)(p'+q'\alpha+r'\alpha^2) = (p+p') + (q+q')\alpha + (r+r')\alpha^2 \in S$, $-(p+q\alpha+r\alpha^2) = (-p) + (-q)\alpha + (-r)\alpha^2 \in S$ and $(p+q\alpha+r\alpha^2)(p'+q'\alpha+r'\alpha^2) = (pp'+2rq') + (pq'+p'q+2rr')\alpha + (pr'+qq'+qr'+rp')\alpha^2 \in S$. By the subring test, $S$ is a ring.

Next, we construct an inverse explicitly to show that $S$ is a subfield. Let $x = p+q\alpha+r\alpha^2 \in S \setminus \{0\}$, where $p,q,r \in \Q$. Then we have $x^2=(p^2+4qr)+(2pq+2r^2)\alpha+(2pr+q^2)\alpha^2$ and $x^3=(p^3+12pqr+2q^3+4r^3)+(3p^2q+6pr^2+6q^2r)\alpha+(3p^2r+3pq^2+6qr^2)\alpha^2$. Let $x^3+k_2x^2+k_1x+k_0=0$. Denote the coefficient of $\alpha^j$ in $x^i$ by $x_{ij}$. Then we have $(x_{30}+x_{31}\alpha+x_{32}\alpha^2)+k_2(x_{20}+x_{21}\alpha+x_{22}\alpha^2)+k_1(p+q\alpha+r\alpha^2)+k_0=0$, that is, $(x_{30}+x_{20}k_2+pk_1+k_0) + (x_{31}+x_{21}k_2+qk_1)\alpha + (x_{32}+x_{22}k_2+rk_1)\alpha^2 = 0$. Since $1,\alpha,\alpha^2$ are linearly independent, comparing coefficients gives us the linear system
\[ \begin{cases}
    k_0 + pk_1 + x_{20}k_2 = -x_{30} \\
    qk_1 + x_{21}k_2 = -x_{31} \\
    rk_1 + x_{22}k_2 = -x_{32}
\end{cases}. \]
We can use the following python code to solve this system.

\begin{minted}[frame=lines, linenos, fontsize=\footnotesize]{python}
import sympy as sp

k0, k1, k2 = sp.symbols('k0 k1 k2')
p, q, r = sp.symbols('p q r')

x30 = p**3 + 12 * p * q * r + 2 * (q**3) + 4 * (r**3)
x31 = 3 * (p**2) * q + 6 * p * (r**2) + 6 * (q**2) * r
x32 = 3 * (p**2) * r + 3 * p * (q**2) + 6 * q * (r**2)
x20 = p**2 + 4 * q * r
x21 = 2 * p * q + 2 * (r**2)
x22 = 2 * p * r + q**2

M = sp.Matrix([
    [1, p, x20, - x30],
    [0, q, x21, - x31],
    [0, r, x22, - x32]
])

solution = sp.solve_linear_system(M, k0, k1, k2)

for var in [k0, k1, k2]:
    print(f"{var} =")
    sp.pprint(sp.simplify(solution[var]))
\end{minted}

The solution is $k_0 = -p^3+6pqr-2q^3-4r^3$, $k_1=3p^2-6qr$ and $k_2=-3p$. Since $x^3+k_2x^2+k_1x+k_0 = 0$, we have $x(x^2+k_2x+k_1+k_0x^{-1}) = 0$. Recall that $\C$ is a field and hence a domain, and $x \neq 0$. Therefore, $x^2+k_2x+k_1+k_0x^{-1} = 0$, that is,
\[x^{-1} = -\frac{1}{k_0}x^2 -\frac{k_2}{k_0}x -\frac{k_1}{k_0}.\]
Since $x$ and $x^2$ are linear combinations of $1,\alpha,\alpha^2$, $x^{-1}$ is also a linear combination of $1,\alpha,\alpha^2$. \epf

\subsection*{Exercise 2}

Assume for contradiction that there exists $a,b \in \Z$ with $b \neq 0$ such that $\lt(\frac{a}{b}\rt)^2=2$. If $a$ is negative and $b$ is positive, then replacing $a$ with $-a$ gives us $\lt(\frac{-a}{b}\rt)^2=\lt(\frac{a}{b}\rt)^2=2$. Therefore, $-a,b$ is also a valid pair of integers satisfying the conditions. If $b$ is negative and $a$ is positive, then replacing $b$ with $-b$ gives us $\lt(\frac{a}{-b}\rt)^2=\lt(\frac{a}{b}\rt)^2=2$. Therefore, $a,-b$ is also a valid pair of integers satisfying the conditions. If both $a$ and $b$ are negative, then replacing $a$ with $-a$ and $b$ with $-b$ gives us $\lt(\frac{-a}{-b}\rt)^2=\lt(\frac{a}{b}\rt)^2=2$. Therefore, $-a,-b$ is also a valid pair of integers satisfying the conditions. Hence, we can assume without loss of generality that $a,b>0$.

Since we assume that $b$ only takes values in $\N$, by the well-ordering principle, there exists a smallest $b$ satisfying the conditions. Note that
\begin{align*}
    \lt(\frac{2b-a}{a-b}\rt)^2 &= \frac{4b^2-4ab+a^2}{a^2-2ab+b^2} \\
    &= \frac{4-4\lt(\frac{a}{b}\rt)+\lt(\frac{a}{b}\rt)^2}{\lt(\frac{a}{b}\rt)^2-2\lt(\frac{a}{b}\rt)+1} \\
    &= \frac{4-4\lt(\frac{a}{b}\rt)+2}{2-2\lt(\frac{a}{b}\rt)+1} \\
    &= \frac{6-4\lt(\frac{a}{b}\rt)}{3-2\lt(\frac{a}{b}\rt)} \\
    &= \frac{6-4\lt(\frac{a}{b}\rt)}{3-2\lt(\frac{a}{b}\rt)} \times \frac{3+2\lt(\frac{a}{b}\rt)}{3+2\lt(\frac{a}{b}\rt)} \\
    &= \frac{18 + 12\lt(\frac{a}{b}\rt) - 12\lt(\frac{a}{b}\rt) - 8\lt(\frac{a}{b}\rt)^2}{9-4\lt(\frac{a}{b}\rt)^2} \\
    &= \frac{18-8\times2}{9-4\times2} \\
    &= 2.
\end{align*}
Recall that $\lt(\frac{a}{b}\rt)^2=2$, so we have $a^2=2b^2$ and thus $2b^2-a^2=0$. Note that
\begin{align*}
    (2b-a)(2b+a) &= 4b^2-a^2 \\
    &= 2b^2 + (2b^2-a^2) \\
    &= 2b^2 \\
    &> 0
\end{align*}
since $b>0$ by assumption. Since $a,b>0$, we have $2b+a>0$. Hence, $2b-a>0$. Also, from $a^2=2b^2$ we have $a^2-b^2=b^2$. Considering $b>0$, we have
\begin{align*}
    b^2 &> 0 \\
    a^2-b^2 &> 0 \\
    (a-b)(a+b) &> 0 \\
    a-b &> 0
\end{align*}
since $a+b>0$. From $2b-a>0$ and $a-b>0$ we know that $2b-a,a-b$ is also a pair of positive integers satisfying the conditions. However, from $2b-a>0$ we have $a<2b$, and hence $a-b<2b-b=b$, contradicting the minimality of $b$. \epf

















\end{document}
